\documentclass[../../main.tex]{subfiles}
\begin{document}

\begin{center}
  {\fontsize{48pt}{58pt}\selectfont \textbf{1968}\normalsize\textbf{年度}}
\end{center}

\begin{table}[htb]
  \centering
  \caption{出題テーマと難易度}
  \begin{tabular}{|c|p{6.2cm}|c|c|c|}
    \hline
    \textbf{問題番号} & \textbf{テーマ} & \textbf{思考量} & \textbf{計算量} & \textbf{総合} \\
    \hline
    第1問 & 命題の真偽と反例・帰納法 & 普 & 易 & 普 \\
    \hline
    第2問 & 三辺の長さを選んで作る三角形の数 & 易 & 易 & 易 \\
    \hline
    第3問 & 平方で閉じる区間の条件 & 普 & 普 & 普 \\
    \hline
    第4問 & 四面体における角の大小関係 & 難 & 普 & 普 \\
    \hline
    第5問 & 放物線上の点から2定点までの距離の平方和 & 易 & 易 & 易 \\
    \hline
    第6問 & 積分で定まる関数の単調性と極限 & 易 & 普 & 普 \\
    \hline
  \end{tabular}
\end{table}

\begin{table}[htb]
  \centering
  \caption{解答}
  \renewcommand{\arraystretch}{1.4}
  \begin{tabular}{|c|c|p{9cm}|}
    \hline
    \textbf{問題番号} & \textbf{小問} & \textbf{答え} \\
    \hline
    第1問 & -- & (1)○，(2)×，(3)○，(4)○，(5)× \\
    \hline
    第2問 & -- & $165$種類 \\
    \hline
    第3問 & (1) & $0\le l\le1$ \\
    \cline{2-3}
    & (2) & $m=1$または$m\le0$ \\
    \cline{2-3}
    & (3) & $S=\{1\}$ \\
    \cline{2-3}
    & (4) & $-\sqrt{l}\le m\le0$ \\
    \hline
    第4問 & -- &  \\
    \hline
    第5問 & -- & $f(x)=2x^4-3x^2-2x+\dfrac{61}{8}$．極小値は$\dfrac{37}{8}$，変曲点は$(-\dfrac{1}{2},8),(\dfrac{1}{2},6)$．概形は解答の図を参照 \\
    \hline
    第6問 & -- & $y$は増加し，$\displaystyle\lim_{x\to1-0}y=\dfrac{3}{4}$ \\
    \hline
  \end{tabular}
\end{table}

\paragraph{試験形態}
6問．試験時間は要確認．

\end{document}
